/section{Beyond-Silicon Materials}
/subsection{Limitations of Silicon}
Silicon has been the base of semiconductor technology because it has good electrical properties, 
reliable manufacturing methods and affordable wafer production.. As transistors keep getting smaller 
more physical and engineering problems appear. Short-channel effects, leakage currents, power density 
and problems, with controlling how transistors work become harder to deal with when the devices get 
very tiny [11]. Because of this scientists are looking for materials that can offer electrical, thermal, 
optical or mechanical qualities.

/subsection{Wide-Bandgap Materials}
Bandgap (WBG) semiconductors have much larger energy bandgaps than silicon. This feature helps devices 
handle electric fields and temperatures. It also allows them to switch faster and might lead to power 
losses [12]. Silicon Carbide (SiC) and Gallium Nitride (GaN) are two WBG materials. They are very useful, 
in power electronics, electric vehicles, renewable-energy systems, frequency electronics and power converters. 
Their material properties help designers create more efficient power systems compared with many conventional 
silicon-based designs [13].

/subsection{2D Materials}
Two‑dimensional materials are made of layers that're extremely thin and can have special electrical and 
optical properties. Graphene, dichalcogenides such as MoS₂ and hexagonal boron nitride are examples of 
two‑dimensional materials that scientists are studying for next‑generation electronics [14]. Some 
two‑dimensional semiconductors have bandgaps and can keep strong electrostatic control even when they 
are very small. Their thickness is one atom thick, which makes them appealing for future transistors 
and, for devices that stack on top of each other. However making these materials on a large scale 
dealing with contact resistance fixing material defects ensuring stability and fitting them into 
the usual semiconductor processes are still big challenges [14] [15].

/subsection{Gallium Nitride / Silicon Carbide}
GaN and SiC are now among the important practical alternatives to silicon for power semiconductor applications. 
I find that GaN devices are especially attractive for frequency and high-speed switching. SiC offers performance, 
in high-voltage and high-temperature environments [13]. I see that these technologies are used in electric 
vehicle power systems, fast chargers, renewable-energy converters, data-center power supplies and industrial 
power electronics. However reliability, manufacturing cost, defect control, packaging and device-specific 
limitations still require research.

/subsection{Future Research}
I think future semiconductor research will not rely on a material to replace silicon. Instead semiconductor 
scientists will mix materials based on what the device needs. WBG materials can help power conversion run 
smoothly while 2D materials can make small transistors, sensors, memory and allow heterogeneous integration. 
Researchers are also looking at ways to join emerging materials with CMOS technologies and advanced 
3D architectures [15]. So the future of semiconductor technology will involve mixes of silicon, 
WBG materials, 2D materials and other emerging materials not a quick and total replacement of silicon.